% Thorium présent dans l'écorce terrestre
Le thorium \ch{_{90}^{237}Th} est radioactif et émetteur $\alpha$. La période
radioactive (ou demi-vie) du thorium 237 est égale à 18 jours. On dispose
initialement ($t=0$) d'une source radioactive de thorium de masse
$m_{0}=\qty{1}{\ug}$.
\begin{enumerate}
    \item Écrire l'équation de sa désintégration sachant qu'elle produit du
          radium \ch{Ra}.

          Calculer la masse de thorium 237 restante aux deux dates
          $t_{1}=\qty{36}{jours}$ et $t_{2}=\qty{6}{mois}$.
    \item À quelles dates la masse de thorium 237 restante est-elle de
          $\qty{0.0156}{\ug}$ et de $\qty{0.0039}{\ug}$~?
    \item À quelle date la masse de thorium 237 restante est-elle voisine de
          $\qty{1}{\ng}$~?
\end{enumerate}

